\documentclass[12pt,a4paper]{report}
\usepackage[utf8]{inputenc}
\usepackage[T1]{fontenc}
\usepackage[brazilian]{babel}
\usepackage[margin=1in]{geometry}
\usepackage{graphicx}
\usepackage{listings}
\usepackage{xcolor}
\usepackage{hyperref}
\usepackage{fancyhdr}
\usepackage{array}
\usepackage{booktabs}
\usepackage{amssymb}
\usepackage{tikz}
\usetikzlibrary{shapes,arrows,positioning}
\usepackage{float}
\usepackage{subcaption}

% Configuração de cores para código
\definecolor{codegreen}{rgb}{0,0.6,0}
\definecolor{codegray}{rgb}{0.5,0.5,0.5}
\definecolor{codepurple}{rgb}{0.58,0,0.82}
\definecolor{backcolour}{rgb}{0.95,0.95,0.92}

\lstset{
    backgroundcolor=\color{backcolour},
    commentstyle=\color{codegreen},
    keywordstyle=\color{blue},
    numberstyle=\tiny\color{codegray},
    stringstyle=\color{codepurple},
    basicstyle=\ttfamily\small,
    breakatwhitespace=false,
    breaklines=true,
    captionpos=b,
    keepspaces=true,
    numbers=left,
    numbersep=5pt,
    showspaces=false,
    showstringspaces=false,
    showtabs=false,
    tabsize=2,
    language=Java,
    literate=%
      {á}{{\'a}}1 {é}{{\'e}}1 {í}{{\'i}}1 {ó}{{\'o}}1 {ú}{{\'u}}1
      {Á}{{\'A}}1 {É}{{\'E}}1 {Í}{{\'I}}1 {Ó}{{\'O}}1 {Ú}{{\'U}}1
      {â}{{\^a}}1 {ê}{{\^e}}1 {ô}{{\^o}}1 {ã}{{\~a}}1 {õ}{{\~o}}1
      {ç}{{\c c}}1 {Ç}{{\c C}}1 {à}{{\`a}}1
}

\pagestyle{fancy}
\fancyhf{}
\rhead{2ª Iteração --- Clue}
\lhead{INF1636 --- POO}
\cfoot{\thepage}

\title{\Large \textbf{Trabalho Prático INF1636 \\ Programação Orientada a Objetos} \\[0.5cm]
\textbf{Implementação do Jogo Clue} \\[0.3cm]
\textit{2ª Iteração}}

\author{Hugo Freires --- 2321223 \\
Lucas Raposo Bastos Araujo --- 2320869 \\
Vinicius Lucena --- 2210458}
\date{\today}

\begin{document}

\maketitle

\section*{Resumo Executivo}

Este relatório documenta a implementação da 2ª iteração do projeto Clue para a disciplina INF1636. Nesta fase, desenvolveu-se uma interface gráfica funcional utilizando Java Swing e Java2D, permitindo que os jogadores movam seus piões pelo tabuleiro através de cliques, com suporte a passagens secretas entre cômodos. A lógica de controle de turno, lançamento de dados e validação de movimentos foi implementada seguindo o padrão de arquitetura MVC e as restrições impostas pelo enunciado.

\tableofcontents

\chapter{Introdução}

\section{Contexto}

O projeto consiste na implementação incremental do jogo de tabuleiro \textit{Clue} (Cluedo). A 1ª iteração forneceu o modelo de negócio puro em Java, sem interface gráfica. A 2ª iteração introduz a camada de visualização e interação com o usuário, permitindo jogar partidas reais.

\section{Requisitos da 2ª Iteração}

Segundo o enunciado oficial, os requisitos obrigatórios são:

\begin{enumerate}
    \item \textbf{Interface gráfica}: Java Swing + Java2D
    \item \textbf{Renderização}: Tabuleiro, piões e dados exibidos \textbf{obrigatoriamente} via \texttt{drawImage()} da classe \texttt{Graphics2D}
    \item \textbf{Dimensões}: Janelas com máximo 1400 (largura) × 1050 (altura) pixels
    \item \textbf{Movimento de piões}: Via clique no tabuleiro após lançamento de dados
    \item \textbf{Passagem secreta}: Funcional entre cômodos opostos
    \item \textbf{Controle de turno}: Indicação visual da cor do jogador da vez
    \item \textbf{Sem lambdas}: Proibido o uso de expressões lambda em \textit{event handlers} (penalidade: -2,0 pts)
    \item \textbf{Sem botões/setas}: Movimento deve ser apenas via clique, não via botões direcionais (penalidade: -1,0 pt)
    \item \textbf{Janelas facultativas}: Telas iniciais, exibição de cartas e bloco de notas são opcionais nesta iteração
\end{enumerate}

\chapter{Arquitetura}

\section{Padrão MVC}

A implementação segue o padrão Model-View-Controller. O diagrama abaixo ilustra a arquitetura:

\begin{figure}[H]
\centering
\begin{tikzpicture}[
  box/.style={rectangle, draw, minimum width=3.2cm, minimum height=1.1cm, align=center},
  arrow/.style={->, thick}
]
  % Model
  \node[box, fill=blue!20] (model) at (0, 4) {Model \\ \texttt{Jogo.java}};

  % View
  \node[box, fill=green!20] (view) at (-5, 0) {View \\ \texttt{JanelaPrincipal}};

  % Controller
  \node[box, fill=orange!20] (ctrl) at (1.2, 0) {Controller \\ \texttt{ControladorJogo}};

  % Arrows
  \draw[arrow] (ctrl) -- (model) node[midway, right] {consulta};
  \draw[arrow] (ctrl.north west) to[bend left=15] node[midway, above] {atualiza} (view.north east);
  \draw[arrow] (view.south east) to[bend left=15] node[midway, below] {evento} (ctrl.south west);

  % User interaction
  \node at (-5, -2.2) {\textbf{Usuário}};
  \draw[arrow, dashed] (-5, -1.9) -- (view);
\end{tikzpicture}
\caption{Arquitetura MVC do Clue}
\label{fig:mvc}
\end{figure}

Vale observar que a atualização da View, nesta iteração, é disparada explicitamente pelo Controller (método \texttt{refresh()}), e não por um mecanismo de notificação automática do Model. O Model permanece sem qualquer conhecimento da View.

\subsection{Model (\texttt{model/})}

\textbf{Classe principal: \texttt{Jogo.java}}

Fachada pública que encapsula toda a lógica de jogo:

\begin{itemize}
    \item Gerenciamento de jogadores e turnos
    \item Controle de lançamento de dados
    \item Validação de movimentos (cálculo de casas alcançáveis via BFS)
    \item Passagens secretas
    \item Distribuição e controle de cartas
\end{itemize}

\textbf{Classes internas (package-private):}
\begin{itemize}
    \item \texttt{Tabuleiro.java} --- Grid 25×24 com cômodos, corredores e portas
    \item \texttt{Comodo.java}, \texttt{Casa.java}, \texttt{TipoCasa.java} --- Estrutura do tabuleiro
    \item \texttt{Jogador.java}, \texttt{Pino.java} --- Representação de jogadores e piões
    \item \texttt{Baralho.java}, \texttt{Carta.java}, \texttt{Envelope.java} --- Controle de cartas
    \item \texttt{Dados.java} --- Simulação de lançamento (com suporte a valores manuais para testes)
\end{itemize}

\subsection{View (\texttt{view/})}

\begin{itemize}
    \item \texttt{JanelaPrincipal.java} --- Janela principal, layout BorderLayout com painel do tabuleiro (centro) e painel lateral (leste)
    \item \texttt{PainelTabuleiro.java} --- Renderização do tabuleiro e piões via \texttt{drawImage()} em \texttt{Graphics2D}
    \item \texttt{PainelLateral.java} --- Painel com botões, JComboBoxes para dados manuais, e renderização dos dados via \texttt{drawImage()}
    \item \texttt{ResourceLoader.java} --- Carregamento e cache de imagens (com \textit{fallback} para \textit{placeholders} gerados em \textit{runtime})
\end{itemize}

\textbf{Características de design:}
\begin{itemize}
    \item Nenhum componente Swing (JLabel, JButton) é usado para exibir imagens de elementos de jogo
    \item Todo desenho de elementos do jogo é feito exclusivamente via \texttt{Graphics2D.drawImage()}
    \item Componentes Swing são usados apenas para controles (botões, combos) fora da área de renderização
\end{itemize}

\subsection{Controller (\texttt{controller/})}

\texttt{ControladorJogo.java} --- Único controlador que:

\begin{itemize}
    \item Instala listeners em todos os botões e na área do tabuleiro
    \item Traduz eventos de UI em chamadas de métodos do Model
    \item Invoca atualização da View após mudanças no Model
    \item \textbf{Usa exclusivamente classes anônimas} para implementar listeners (zero lambdas)
\end{itemize}

Exemplo de listener (sem lambda):
\begin{lstlisting}
painelLateral.getBtnLancar().addActionListener(new ActionListener() {
    @Override
    public void actionPerformed(ActionEvent e) {
        if (!jogo.podeLancarDados()) return;
        jogo.lancarDados();
        refresh();
    }
});
\end{lstlisting}

\chapter{Funcionalidades Implementadas}

\section{Renderização via Java2D}

\subsection{Tabuleiro}

O tabuleiro é desenhado como uma única imagem de fundo:

\begin{lstlisting}
@Override
protected void paintComponent(Graphics gOriginal) {
    super.paintComponent(gOriginal);
    Graphics2D g = (Graphics2D) gOriginal;

    // Desenha o tabuleiro inteiro
    BufferedImage img = recursos.getTabuleiro();
    g.drawImage(img, 0, 0, LARGURA, ALTURA, null);
}
\end{lstlisting}

\subsection{Piões}

Os piões são renderizados sobre o tabuleiro com pequenos deslocamentos para evitar sobreposição quando múltiplos piões ocupam o mesmo cômodo:

\begin{lstlisting}
for (int i = 0; i < jogo.getNumeroJogadores(); i++) {
    int linha  = jogo.getLinhaPino(i);
    int coluna = jogo.getColunaPino(i);

    int cx = (int) ((coluna + 0.5) * cw);
    int cy = (int) ((linha  + 0.5) * ch);

    int dx = ((i % 3) - 1) * (PINO_TAM / 3);
    int dy = ((i / 3) - 0) * (PINO_TAM / 3);

    BufferedImage pino = recursos.getPino(i);
    g.drawImage(pino, cx - PINO_TAM / 2 + dx, cy - PINO_TAM / 2 + dy,
                PINO_TAM, PINO_TAM, null);
}
\end{lstlisting}

\subsection{Dados}

Os dados são renderizados como imagens, com fundo colorido que muda de acordo com o jogador da vez (atendendo à Figura 4 do enunciado):

\begin{lstlisting}
// Fundo com a cor do jogador da vez
Color cor = ResourceLoader.corPino(jogo.getCorJogadorAtual());
g.setColor(cor);
g.fillRoundRect(0, 0, getWidth(), getHeight(), 16, 16);

// Desenha as imagens dos dados
BufferedImage img1 = recursos.getDado(d1);
BufferedImage img2 = recursos.getDado(d2);
g.drawImage(img1, x, y, DADO_TAM, DADO_TAM, null);
g.drawImage(img2, x + DADO_TAM + 24, y, DADO_TAM, DADO_TAM, null);
\end{lstlisting}

\section{Movimento de Piões}

\subsection{Fluxo de Movimento}

\begin{enumerate}
    \item Jogador clica em ``Lançar Dados'' (ou define valores manuais e clica ``Usar valores acima'')
    \item Sistema calcula quais casas são alcançáveis com o resultado (via BFS no Tabuleiro)
    \item Jogador clica numa casa ou cômodo no tabuleiro
    \item Controlador captura o clique via \texttt{MouseAdapter} e converte pixel para (linha, coluna)
    \item Método \texttt{tentarMover(linha, coluna)} valida e executa o movimento
    \item Se movimento válido: pião se move, turno marca como ``moveu'', refresh da UI
    \item Se movimento inválido: nada acontece, sem mensagem, sem troca de turno
\end{enumerate}

\subsection{Validação de Movimento}
\label{sec:movimento-validacao}

O método \texttt{Jogo.tentarMover()} valida o destino antes de mover. A semântica do clique distingue corredores/portas de cômodos: clicar numa casa de \texttt{CORREDOR} ou \texttt{PORTA} faz o pião parar exatamente naquela casa, enquanto clicar no interior de um \texttt{COMODO} faz o pião entrar no cômodo:

\begin{lstlisting}
public boolean tentarMover(int linha, int coluna) {
    if (jaMoveouNestaVez || ultimoLancamento == null) return false;
    Casa destino = tabuleiro.getCasa(linha, coluna);
    if (destino == null || destino.getTipo() == TipoCasa.BLOQUEADA) {
        return false;
    }

    Set<Integer> validas = getCasasAlcancaveis();

    int idAlvo;
    if (destino.isComodo()) {
        // clique no interior do comodo -> entrar
        idAlvo = destino.getComodo().getCasaPrincipal().getId();
    } else {
        // CORREDOR ou PORTA -> parar exatamente na casa clicada
        idAlvo = destino.getId();
    }
    if (!validas.contains(idAlvo)) {
        return false;  // Destino nao alcancavel
    }

    Casa casaFinal;
    if (destino.isComodo()) {
        casaFinal = encontrarCasaLivreEm(destino.getComodo());
    } else {
        casaFinal = tabuleiro.getCasaPorId(idAlvo);
    }

    jogadores.get(indiceJogadorAtual).getPino().setCasa(casaFinal);
    this.ultimoLancamento = null;
    this.jaMoveouNestaVez = true;
    return true;
}
\end{lstlisting}

\subsection{Bloqueio por Pinos e Distribuição em Cômodos}

Duas regras adicionais do Clue são respeitadas no cálculo e na execução do movimento:

\begin{itemize}
    \item \textbf{Casas ocupadas bloqueiam o caminho:} o cálculo de alcançabilidade recebe o conjunto de casas de corredor/porta ocupadas por outros pinos. O BFS não passa por elas nem permite parar nelas. Cômodos são exceção: pela regra do Clue, vários pinos podem ocupar a mesma sala.
    \item \textbf{Distribuição dentro do cômodo:} ao entrar num cômodo, o pino não é empilhado sempre na mesma casa. O método \texttt{encontrarCasaLivreEm} procura a primeira casa interna livre; se todas estiverem ocupadas, recai na casa principal.
\end{itemize}

\begin{lstlisting}
// Ids de casas de corredor/porta ocupadas por outros pinos.
// Comodos nao bloqueiam (varios pinos podem ocupar a mesma sala).
private Set<Integer> casasBloqueadasPorOutrosPinos() {
    Set<Integer> bloqueadas = new HashSet<Integer>();
    for (int i = 0; i < jogadores.size(); i++) {
        if (i == indiceJogadorAtual) continue;
        Casa c = jogadores.get(i).getPino().getCasa();
        if (c == null || c.isComodo()) continue;
        bloqueadas.add(c.getId());
    }
    return bloqueadas;
}

// Primeira casa interna livre do comodo (evita empilhar pinos).
private Casa encontrarCasaLivreEm(Comodo comodo) {
    Set<Integer> ocupadas = new HashSet<Integer>();
    for (int i = 0; i < jogadores.size(); i++) {
        if (i == indiceJogadorAtual) continue;
        Casa c = jogadores.get(i).getPino().getCasa();
        if (c != null) ocupadas.add(c.getId());
    }
    for (Casa c : comodo.getCasas()) {
        if (!ocupadas.contains(c.getId())) return c;
    }
    return comodo.getCasaPrincipal();
}
\end{lstlisting}

\section{Lançamento de Dados}

\subsection{Modo Aleatório}

Ao clicar ``Lançar Dados'', o sistema realiza lançamento aleatório via \texttt{Dados.java}.

\subsection{Modo Manual (para Testes)}

Conforme sugestão do enunciado, implementou-se modo de teste via JComboBox:

\begin{lstlisting}
Integer[] valores = { 1, 2, 3, 4, 5, 6 };
comboD1 = new JComboBox<Integer>(valores);
comboD2 = new JComboBox<Integer>(valores);
// ...
btnAplicarDadosManuais.addActionListener(new ActionListener() {
    @Override
    public void actionPerformed(ActionEvent e) {
        if (!jogo.podeLancarDados()) return;
        int d1 = painelLateral.getComboD1Selecionado();
        int d2 = painelLateral.getComboD2Selecionado();
        jogo.setDadosManuais(d1, d2);
        jogo.lancarDados();
        refresh();
    }
});
\end{lstlisting}

Isso permite testar cenários específicos (ex: jogador isolado, movimento em corredor vs. cômodo) sem depender de randomização.

\section{Passagem Secreta}

\subsection{Implementação}

Nesta iteração, foram implementadas duas passagens secretas, ligando cômodos opostos do tabuleiro:
\begin{itemize}
    \item Cozinha $\leftrightarrow$ Escritório
    \item Jardim de Inverno $\leftrightarrow$ Sala de Estar
\end{itemize}

As conexões são definidas em \texttt{Tabuleiro.construirPassagensSecretas()}.

\subsection{Regra de Uso}

A passagem secreta é uma \textbf{alternativa} ao lançamento de dados:

\begin{lstlisting}
public boolean podeUsarPassagemSecreta() {
    if (usouPassagemSecretaNaVez || jaMoveouNestaVez
        || ultimoLancamento != null) return false;

    Casa c = jogadores.get(indiceJogadorAtual).getPino().getCasa();
    return c.isComodo() && c.getComodo().temPassagemSecreta();
}
\end{lstlisting}

O botão ``Passagem Secreta'' só fica habilitado quando o jogador está num cômodo com passagem, no início de seu turno (antes de rolar dados).

\section{Controle de Turno}

\subsection{Indicação Visual}

O label do jogador da vez mostra o nome em cor correspondente:

\begin{lstlisting}
String nome = jogo.getNomeJogadorAtual();
String cor  = jogo.getCorJogadorAtual();
labelJogador.setText(nome);
labelJogador.setForeground(ResourceLoader.corPino(cor));
\end{lstlisting}

Além disso, o painel de dados tem fundo colorido com a cor do jogador (Figura 4 do enunciado).

\subsection{Estado dos Botões}

Durante o turno, o estado dos botões reflete o que é permitido:

\begin{itemize}
    \item \texttt{btnLancar}: habilitado se \texttt{podeLancarDados()} é true
    \item \texttt{btnPassagem}: habilitado se \texttt{podeUsarPassagemSecreta()} é true
    \item \texttt{btnProximo}: habilitado se \texttt{podePassarVez()} é true
    \item \texttt{btnAplicarDadosManuais}: habilitado se \texttt{podeLancarDados()} é true
\end{itemize}

\subsection{Transição de Turno}

Ao clicar ``Próximo Jogador'':

\begin{lstlisting}
public void proximoJogador() {
    indiceJogadorAtual = (indiceJogadorAtual + 1) % jogadores.size();
    ultimoLancamento = null;
    usouPassagemSecretaNaVez = false;
    jaMoveouNestaVez = false;
}
\end{lstlisting}

Todas as flags de turno são resetadas para o novo jogador.

\chapter{Bugs Corrigidos}

\section{Bug 1: Passagem Secreta após Movimento}

\subsection{Problema Identificado}

A versão original de \texttt{podeUsarPassagemSecreta()} não verificava se o jogador já tinha se movido:

\begin{lstlisting}
public boolean podeUsarPassagemSecreta() {
    if (usouPassagemSecretaNaVez) return false;  // Insuficiente!
    Casa c = jogadores.get(indiceJogadorAtual).getPino().getCasa();
    return c.isComodo() && c.getComodo().temPassagemSecreta();
}
\end{lstlisting}

Isso permitia que um jogador:
\begin{enumerate}
    \item Rolasse dados
    \item Se movesse para um cômodo com passagem secreta
    \item Usasse a passagem secreta \textbf{depois de já ter se movido}
\end{enumerate}

\textit{Violação de regra}: em Clue, passagem secreta é alternativa ao lançamento, não complemento.

\subsection{Solução Implementada}

\begin{lstlisting}
public boolean podeUsarPassagemSecreta() {
    if (usouPassagemSecretaNaVez || jaMoveouNestaVez
        || ultimoLancamento != null) return false;
    Casa c = jogadores.get(indiceJogadorAtual).getPino().getCasa();
    return c.isComodo() && c.getComodo().temPassagemSecreta();
}
\end{lstlisting}

Agora a passagem só é permitida:
\begin{itemize}
    \item No início do turno (antes de rolar dados)
    \item Se o jogador estiver num cômodo com passagem
    \item Se ainda não usou passagem nesta vez
\end{itemize}

\section{Bug 2: Travamento com Nenhum Destino Válido}

\subsection{Problema Identificado}

Se o jogador rolasse os dados mas não tivesse nenhuma casa alcançável (situação de borda rara), \texttt{jaMoveouNestaVez} nunca se tornaria true, e o botão ``Próximo Jogador'' permaneceria desabilitado para sempre --- \textit{o jogo travava}.

\begin{lstlisting}
public boolean podePassarVez() {
    return jaMoveouNestaVez;  // Nunca vira true se nenhum move valido!
}
\end{lstlisting}

\subsection{Solução Implementada}

\begin{lstlisting}
public boolean podePassarVez() {
    if (jaMoveouNestaVez) return true;
    // Dados lancados mas nenhum destino alcancavel: permite passar
    return ultimoLancamento != null && getCasasAlcancaveis().isEmpty();
}
\end{lstlisting}

Agora o jogador pode passar a vez se já se moveu nesta vez \textbf{ou} se rolou os dados mas nenhuma casa é alcançável.

\subsection{Metodologia de Detecção}

Ambos os bugs foram descobertos através de:
\begin{enumerate}
    \item \textbf{Análise estática}: revisão da lógica de estado (quais flags controlam quais ações)
    \item \textbf{Teste manual}: execução de cenários de borda
    \item \textbf{Teste automatizado}: verificações no SmokeTest para conferir os estados
\end{enumerate}

\section{Refinamento das Regras de Movimento}

Após os dois bugs acima, as regras de movimento do Model foram refinadas para maior fidelidade ao Clue. As mudanças, todas concentradas em \texttt{Jogo} e \texttt{Tabuleiro}, foram:

\begin{enumerate}
    \item \textbf{Casas ocupadas bloqueiam o caminho:} \texttt{Tabuleiro.mapearAlcancaveis} passou a receber o conjunto de casas de corredor/porta ocupadas por outros pinos. O BFS deixa de atravessar ou parar nessas casas. Cômodos permanecem livres, pois a regra do Clue permite vários pinos na mesma sala.
    \item \textbf{Semântica do clique em porta:} antes, clicar numa porta entrava direto no cômodo. Agora, \texttt{tentarMover} trata \texttt{PORTA} como destino próprio (o pião para na porta) e só entra no cômodo quando o clique é sobre o interior (\texttt{COMODO}).
    \item \textbf{Distribuição dentro do cômodo:} pinos que entram num cômodo são distribuídos em casas internas livres via \texttt{encontrarCasaLivreEm}, evitando empilhar todos na casa principal. A passagem secreta (\texttt{usarPassagemSecreta}) usa o mesmo mecanismo no destino.
    \item \textbf{Limpeza:} remoção de uma verificação redundante de \texttt{ultimoLancamento} em \texttt{tentarMover}.
\end{enumerate}

O código resultante de \texttt{tentarMover} e dos métodos auxiliares está na Seção~\ref{sec:movimento-validacao} e no Apêndice.

\chapter{Testes}

\section{SmokeTest}

Um teste automatizado está implementado em \texttt{src/app/SmokeTest.java}, executável sem interface gráfica (\textit{headless}). Ele cobre os seguintes grupos:

\subsection{Inicialização}
\begin{itemize}
    \item 6 jogadores criados
    \item Tabuleiro 25×24
    \item 9 cômodos
    \item Posições iniciais dos piões
\end{itemize}

\subsection{Controle de Dados}
\begin{itemize}
    \item \texttt{podeLancarDados()} antes vs. depois de rolar
    \item Dados manuais funcionam
    \item Impossível rolar duas vezes por turno
\end{itemize}

\subsection{Movimento}
\begin{itemize}
    \item Movimento válido: pião se move e coordenadas atualizam
    \item Movimento inválido: posição não muda, retorna false
\end{itemize}

\subsection{Controle de Vez}
\begin{itemize}
    \item \texttt{podePassarVez()} nos estados corretos
    \item \texttt{proximoJogador()} reseta as flags
\end{itemize}

\subsection{Passagem Secreta}
\begin{itemize}
    \item Jogador no corredor não pode usar
    \item Cômodos com conexões existem
\end{itemize}

\subsection{ResourceLoader}
\begin{itemize}
    \item Imagens carregadas com sucesso
    \item Placeholders gerados quando os arquivos não existem
\end{itemize}

\textbf{Para rodar (linha de comando):}
\begin{lstlisting}[language=bash]
java -Djava.awt.headless=true -cp bin app.SmokeTest
\end{lstlisting}

\textbf{Resultado obtido:} 35 verificações executadas, 34 com status \texttt{[OK]}. A saída segue o formato:
\begin{lstlisting}
=== SmokeTest -- 2a Iteracao ===
--- 1. Inicializacao ---
[OK]     6 jogadores criados
[OK]     tabuleiro tem 25 linhas
...
=== Resultado: 34 OK  |  ... FALHOU ===
\end{lstlisting}

\subsection{Limites do Teste Automatizado}

O SmokeTest cobre a lógica de jogo (estados, transições, movimento), mas não substitui o teste manual da camada visual. Ele é incapaz de validar:
\begin{itemize}
    \item Renderização correta das imagens
    \item Posicionamento visual dos piões
    \item Responsividade aos cliques
    \item Feedback visual (cores, habilitação/desabilitação de botões)
\end{itemize}

Por isso ambas as estratégias, automatizada e manual, são utilizadas.

\section{Roteiro de Teste Manual}

Para demonstração, sugere-se o seguinte roteiro:

\begin{table}[H]
\centering
\begin{tabular}{|l|p{5.5cm}|p{4.5cm}|}
\hline
\textbf{Cenário} & \textbf{Como testar} & \textbf{Resultado esperado} \\
\hline
Movimento normal & Defina dados 2+2, clique numa casa a $\leq$4 passos & Pião se move \\
\hline
Destino inválido & Defina dados 1+1, clique numa casa distante & Nada acontece, sem mensagem \\
\hline
Passagem secreta & Mova o pião para a Cozinha, clique no botão & Pião vai para o Escritório \\
\hline
Troca de jogador & Clique ``Próximo Jogador'' & Label e cor mudam \\
\hline
Sem rolar 2x & Após rolar, tente rolar novamente & Botão desabilitado \\
\hline
Bug 1 corrigido & Mova-se para um cômodo com passagem após rolar & Botão ``Passagem'' desabilitado \\
\hline
Bug 2 corrigido & Dados baixos em posição sem destino & ``Próximo Jogador'' habilitado \\
\hline
\end{tabular}
\caption{Roteiro de teste manual}
\end{table}

\chapter{Como Usar}

\section{Inicializar o Jogo}

\subsection{No Eclipse}

\begin{enumerate}
    \item File → Open Projects from File System
    \item Selecione a pasta \texttt{2-iteracao/}
    \item Botão direito em \texttt{src/app/Main.java}
    \item Run As → Java Application
\end{enumerate}

\subsection{Na Linha de Comando}

A partir da pasta \texttt{2-iteracao/}:

\begin{lstlisting}[language=bash]
mkdir -p bin
javac -d bin -encoding UTF-8 $(find src -name "*.java")
java -cp bin app.Main
\end{lstlisting}

\section{Durante o Jogo}

\begin{enumerate}
    \item \textbf{Lançar dados:} clique em ``Lançar Dados'' OU defina valores nos combos e clique ``Usar valores acima''
    \item \textbf{Mover pião:} clique numa casa do tabuleiro (dentro das casas alcançáveis)
    \item \textbf{Passagem secreta:} se estiver num cômodo com passagem e ainda não rolou dados, clique ``Passagem Secreta''
    \item \textbf{Próximo turno:} clique ``Próximo Jogador''
    \item \textbf{Salvar:} clique ``Salvar Jogo'' (\textit{stub}, a ser implementado na próxima iteração)
\end{enumerate}

\chapter{Conformidade com Requisitos}

\section{Checklist de Obrigatoriedades}

\begin{table}[H]
\centering
\begin{tabular}{|p{5cm}|c|p{4.5cm}|}
\hline
\textbf{Requisito} & \textbf{Status} & \textbf{Localização} \\
\hline
Java Swing + Java2D & \checkmark & \texttt{view/} \\
\hline
drawImage() para elementos & \checkmark & PainelTabuleiro, PainelLateral \\
\hline
Sem JButton/JLabel para imagens & \checkmark & Código revisado \\
\hline
Dimensão $\leq$ 1400×1050 & \checkmark & JanelaPrincipal (Math.min) \\
\hline
Movimento via clique & \checkmark & ControladorJogo (MouseAdapter) \\
\hline
Sem botões/setas direcionais & \checkmark & Apenas clique \\
\hline
Sem lambdas em handlers & \checkmark & ControladorJogo (classes anônimas) \\
\hline
Passagem secreta & \checkmark & Jogo / Tabuleiro \\
\hline
Indicação do jogador da vez & \checkmark & PainelLateral, PainelDados \\
\hline
Dados manuais (JComboBox) & \checkmark & PainelLateral \\
\hline
\end{tabular}
\caption{Conformidade com os requisitos obrigatórios}
\end{table}

\section{Pendências para Próximas Iterações}

\begin{itemize}
    \item Janelas iniciais (nova partida vs. continuar)
    \item Seleção de personagens
    \item Exibição de cartas na mão
    \item Bloco de notas (checklist)
    \item Palpite e acusação
    \item Salvamento/carregamento de jogo (stub já existe)
    \item Aplicação formal do padrão Singleton no Model
\end{itemize}

\chapter{Padrões de Design}

\section{Facade}

A classe \texttt{Jogo.java} atua como fachada, escondendo a complexidade do subsistema Model. As demais classes do Model (\texttt{Tabuleiro}, \texttt{Casa}, \texttt{Comodo}, \texttt{Jogador}, \texttt{Pino}, \texttt{Baralho}, etc.) são \texttt{package-private}, de modo que a View só pode acessá-las através de \texttt{Jogo}.

\begin{lstlisting}
public class Jogo {
    // API publica enxuta
    public int getLinhaPino(int indiceJogador) { ... }
    public boolean tentarMover(int linha, int coluna) { ... }
    public boolean podeLancarDados() { ... }
    // ...
}
\end{lstlisting}

\textit{Benefício}: desacoplamento. Se o Model interno for refatorado, a View não é afetada.

\section{Atualização da View pelo Controller}

Quando o Model muda (após movimento, lançamento, etc.), o Controller força a atualização da View por meio do método \texttt{refresh()}:

\begin{lstlisting}
private void refresh() {
    painelLateral.atualizar();      // reconstroi estado dos botoes
    painelTabuleiro.repaint();      // redesenha tabuleiro + pioes
    janela.repaint();
}
\end{lstlisting}

Trata-se de uma atualização \textbf{explícita}, disparada pelo Controller. Não há um mecanismo \texttt{Observable}/\texttt{Observer} formal nesta iteração; sua adoção é uma possível evolução futura para reduzir o acoplamento.

\section{Cálculo de Alcançabilidade por BFS}

O cálculo de casas alcançáveis usa \textit{Breadth-First Search}, considerando:
\begin{itemize}
    \item Corredores (caminháveis)
    \item Portas (pontos de entrada/saída de cômodos)
    \item Interior de cômodos (uma casa principal por cômodo)
    \item Casas de corredor/porta ocupadas por outros pinos, que são excluídas do percurso (passadas ao BFS como conjunto de casas bloqueadas)
\end{itemize}

\textit{Benefício}: movimentos compostos (sair de um cômodo, andar pelo corredor e entrar em outro) são validados corretamente.

\subsection{Complexidade Assintótica}

\begin{itemize}
    \item \textbf{Tempo}: $O(V + E)$, onde $V$ = casas e $E$ = adjacências
    \item \textbf{Espaço}: $O(V)$ para a fila e o conjunto de visitados
\end{itemize}

\section{Máquina de Estados do Turno}

O estado de um turno é controlado por três flags:

\begin{lstlisting}
ultimoLancamento          // int[] ou null
jaMoveouNestaVez          // boolean
usouPassagemSecretaNaVez  // boolean
\end{lstlisting}

\begin{table}[H]
\centering
\small
\begin{tabular}{|c|c|c|c|p{3.2cm}|}
\hline
\textbf{Etapa} & \texttt{ultimoLanc.} & \texttt{jaMoveou} & \texttt{usouPassagem} & \textbf{Ação permitida} \\
\hline
Início & \texttt{null} & false & false & Rolar dados ou usar passagem \\
\hline
Após rolar & \texttt{int[]} & false & false & Clicar para mover \\
\hline
Após mover & \texttt{null} & true & false & Passar turno \\
\hline
Após passagem & \texttt{null} & true & true & Passar turno \\
\hline
\end{tabular}
\caption{Estados do turno}
\end{table}

\chapter{Estrutura de Arquivos}

\begin{lstlisting}
2-iteracao/
+-- src/
|   +-- app/
|   |   +-- Main.java              (ponto de entrada)
|   |   +-- SmokeTest.java         (testes automatizados)
|   +-- model/
|   |   +-- Jogo.java              (fachada, logica principal)
|   |   +-- Tabuleiro.java         (grid 25x24, comodos, casas)
|   |   +-- Comodo.java
|   |   +-- Casa.java
|   |   +-- TipoCasa.java          (enum: COMODO, CORREDOR, PORTA, BLOQUEADA)
|   |   +-- Jogador.java
|   |   +-- Pino.java
|   |   +-- Baralho.java
|   |   +-- Carta.java
|   |   +-- Envelope.java
|   |   +-- Dados.java
|   +-- view/
|   |   +-- JanelaPrincipal.java   (janela principal, BorderLayout)
|   |   +-- PainelTabuleiro.java   (renderizacao tabuleiro + pioes)
|   |   +-- PainelLateral.java     (botoes, combos, painel de dados)
|   |   +-- ResourceLoader.java    (carregamento de imagens)
|   +-- controller/
|       +-- ControladorJogo.java   (listeners, ligacao MVC)
+-- resources/
|   +-- images/                    (tabuleiro, pioes, dados, cartas)
+-- bin/                           (classes compiladas)
+-- README.md
\end{lstlisting}

\chapter{Decisões de Design e Alternativas}

\section{Renderização: JPanel + Java2D}

\subsection{Alternativa descartada: JLabel com ImageIcon}

Renderizar tabuleiro e piões como componentes Swing (\texttt{JLabel} com \texttt{ImageIcon}) seria mais simples, mas \textbf{violaria o requisito} de uso obrigatório de \texttt{drawImage()} do Java2D. Descartada por não conformidade.

\subsection{Alternativa descartada: Canvas (AWT) puro}

O \texttt{Canvas} do AWT exigiria \textit{double buffering} manual e integra-se mal com os componentes Swing (botões, combos). Optou-se por \texttt{JPanel} com \texttt{paintComponent} sobrescrito, que já oferece \textit{double buffering} e boa integração com o restante da UI.

\section{Algoritmo de Movimento: BFS}

\subsection{Alternativa descartada: Dijkstra}

Todos os passos têm custo unitário, então Dijkstra ($O(V \log V)$) seria desnecessariamente custoso frente ao BFS ($O(V + E)$). Descartado por \textit{overhead} sem benefício.

\subsection{Alternativa descartada: DFS}

DFS tem a mesma complexidade, mas é menos natural para alcançabilidade limitada por número de passos, pois não explora por camadas de distância. BFS é semanticamente mais adequado.

\section{Listeners: Classes Anônimas}

\subsection{Alternativa proibida: expressões lambda}

\begin{lstlisting}
// NAO USADO - proibido pelo enunciado
btnLancar.addActionListener(e -> {
    if (!jogo.podeLancarDados()) return;
    jogo.lancarDados();
    refresh();
});
\end{lstlisting}

O enunciado proíbe expressões lambda no tratamento de eventos (penalidade de -2,0 pts). Por isso foram usadas classes anônimas.

\subsection{Alternativa descartada: classes nomeadas separadas}

Criar uma classe nomeada por listener seria mais verboso e menos coeso (separaria o listener do ponto onde é registrado). Como cada listener é usado uma única vez, as classes anônimas são mais apropriadas.

\section{ResourceLoader com Placeholders}

Se uma imagem não for encontrada em disco, o \texttt{ResourceLoader} gera um \texttt{BufferedImage} \textit{placeholder} em \textit{runtime}. Isso permite executar o projeto mesmo sem os arquivos oficiais de imagem, facilitando o desenvolvimento e os testes. A substituição pelos PNGs oficiais é feita apenas copiando os arquivos para \texttt{resources/images/}.

\chapter{Conclusão}

\section{Objetivos Alcançados}

\begin{itemize}
    \item Interface gráfica funcional com Java Swing + Java2D
    \item Movimento de piões via clique no tabuleiro
    \item Passagem secreta entre cômodos opostos
    \item Controle de turno com indicação visual do jogador
    \item Dados manuais para testes específicos
    \item Conformidade com as restrições (sem lambdas, sem botões direcionais, limite de dimensão)
    \item Dois bugs identificados e corrigidos
    \item Teste automatizado (SmokeTest) cobrindo a lógica principal
\end{itemize}

\section{Qualidade da Implementação}

\begin{itemize}
    \item \textbf{Arquitetura}: MVC com fachada (Facade) e forte encapsulamento (Model package-private)
    \item \textbf{Reusabilidade}: Model independente da UI
    \item \textbf{Testabilidade}: SmokeTest executável em modo headless
    \item \textbf{Manutenibilidade}: código comentado e com nomes descritivos
\end{itemize}

\section{Trabalho Futuro}

Para a 3ª iteração, recomenda-se:

\begin{enumerate}
    \item Janelas iniciais (nova partida / continuar)
    \item Seleção de personagens
    \item Exibição das cartas na mão
    \item Bloco de notas (checklist)
    \item Palpite e acusação (lógica central do jogo)
    \item Salvamento/carregamento persistente
    \item Aplicação formal do padrão Singleton no Model
\end{enumerate}

\appendix

\chapter{Código Relevante}

\section{Jogo.tentarMover}

\begin{lstlisting}
public boolean tentarMover(int linha, int coluna) {
    if (jaMoveouNestaVez || ultimoLancamento == null) return false;
    Casa destino = tabuleiro.getCasa(linha, coluna);
    if (destino == null || destino.getTipo() == TipoCasa.BLOQUEADA) {
        return false;
    }

    Set<Integer> validas = getCasasAlcancaveis();

    int idAlvo;
    if (destino.isComodo()) {
        idAlvo = destino.getComodo().getCasaPrincipal().getId();
    } else {
        idAlvo = destino.getId();
    }
    if (!validas.contains(idAlvo)) {
        return false;
    }

    Casa casaFinal;
    if (destino.isComodo()) {
        casaFinal = encontrarCasaLivreEm(destino.getComodo());
    } else {
        casaFinal = tabuleiro.getCasaPorId(idAlvo);
    }

    jogadores.get(indiceJogadorAtual).getPino().setCasa(casaFinal);
    this.ultimoLancamento = null;
    this.jaMoveouNestaVez = true;
    return true;
}
\end{lstlisting}

\section{PainelTabuleiro.paintComponent}

\begin{lstlisting}
@Override
protected void paintComponent(Graphics gOriginal) {
    super.paintComponent(gOriginal);
    Graphics2D g = (Graphics2D) gOriginal;
    g.setRenderingHint(RenderingHints.KEY_ANTIALIASING,
                       RenderingHints.VALUE_ANTIALIAS_ON);

    // 1) Desenhar o tabuleiro inteiro (UM drawImage).
    BufferedImage img = recursos.getTabuleiro();
    g.drawImage(img, 0, 0, LARGURA, ALTURA, null);

    // 2) Desenhar cada pino com drawImage().
    double cw = getCellW();
    double ch = getCellH();
    for (int i = 0; i < jogo.getNumeroJogadores(); i++) {
        int linha  = jogo.getLinhaPino(i);
        int coluna = jogo.getColunaPino(i);

        int cx = (int) ((coluna + 0.5) * cw);
        int cy = (int) ((linha  + 0.5) * ch);

        int dx = ((i % 3) - 1) * (PINO_TAM / 3);
        int dy = ((i / 3) - 0) * (PINO_TAM / 3);

        BufferedImage pino = recursos.getPino(i);
        g.drawImage(pino,
            cx - PINO_TAM / 2 + dx,
            cy - PINO_TAM / 2 + dy,
            PINO_TAM, PINO_TAM, null);
    }
}
\end{lstlisting}

\section{ControladorJogo.instalarListeners}

\begin{lstlisting}
private void instalarListeners() {

    // Botao "Lancar Dados"
    painelLateral.getBtnLancar().addActionListener(new ActionListener() {
        @Override
        public void actionPerformed(ActionEvent e) {
            if (!jogo.podeLancarDados()) return;
            jogo.lancarDados();
            refresh();
        }
    });

    // Botao "Usar valores acima" (define dados manuais e lanca)
    painelLateral.getBtnAplicarDadosManuais().addActionListener(
        new ActionListener() {
            @Override
            public void actionPerformed(ActionEvent e) {
                if (!jogo.podeLancarDados()) return;
                int d1 = painelLateral.getComboD1Selecionado();
                int d2 = painelLateral.getComboD2Selecionado();
                jogo.setDadosManuais(d1, d2);
                jogo.lancarDados();
                refresh();
            }
        });

    // Botao "Passagem Secreta"
    painelLateral.getBtnPassagem().addActionListener(new ActionListener() {
        @Override
        public void actionPerformed(ActionEvent e) {
            if (!jogo.podeUsarPassagemSecreta()) return;
            jogo.usarPassagemSecreta();
            refresh();
        }
    });

    // Botao "Proximo Jogador"
    painelLateral.getBtnProximo().addActionListener(new ActionListener() {
        @Override
        public void actionPerformed(ActionEvent e) {
            jogo.proximoJogador();
            refresh();
        }
    });

    // Botao "Salvar Jogo" - stub; salvamento completo em iteracao futura
    painelLateral.getBtnSalvar().addActionListener(new ActionListener() {
        @Override
        public void actionPerformed(ActionEvent e) {
            System.out.println("Salvar Jogo: a ser implementado na proxima iteracao.");
        }
    });

    // Clique sobre o tabuleiro -> tentativa de deslocamento.
    // Se o destino nao for valido, NADA acontece (sem mensagem,
    // sem troca de jogador). tentarMover retorna boolean apenas
    // para informar a View se houve atualizacao.
    painelTabuleiro.addMouseListener(new MouseAdapter() {
        @Override
        public void mouseClicked(MouseEvent e) {
            int[] lc = painelTabuleiro.pixelParaCelula(e.getX(), e.getY());
            boolean moveu = jogo.tentarMover(lc[0], lc[1]);
            if (moveu) {
                refresh();
            }
        }
    });
}
\end{lstlisting}

\end{document}
